\BourbakiExerciseGroup

\begin{enumerate}
\def\labelenumi{\arabic{enumi})}
\item
  Let \(f\) be a continuous mapping of an open interval \(\mathbf{I} \subset \mathbb{R}\) into \(\mathbb{R}\) . Show that if \(f(\mathbf{I})\) is open and if, for each \(y \in \mathbb{R}\) , the set \(\overline{f}^1(y)\) has at most two distinct points, then \(f\) is monotone.
\item
  Let B be a base of R considered as a vector space over the field Q (``Hamel base''). B is uncountable (§ 2, Exercise 2). Let φ be a bijection of a subset C ≠ B of B onto the set B. Define a mapping f of R into R as follows: if \(x = \sum_{\xi \in B} \lambda(\xi) \xi\) , where \(\lambda(\xi) \in \mathbf{Q}\) and \(\lambda(\xi) = 0\) for all but a finite number of elements \(\xi \in B\) , then \(f(x) = \sum_{\xi \in C} \lambda(\xi) \varphi(\xi)\) . Show that \(f(x + y) = f(x) + f(y)\) , but that, for each \(z \in R\) , \(f(z)\) is a dense subset of R, which implies that f is not bounded above nor below in any interval of R.
\item
  \begin{enumerate}
  \def\labelenumii{\alph{enumii})}
  \tightlist
  \item
    For each interval \(\mathbf{I} \subset \mathbb{R}\) , let \(\mathbf{G}(\mathbf{I})\) denote the group of all homeomorphisms of \(\mathbf{I}\) onto itself. Show that if \(\mathbf{I}\) and \(\mathbf{J}\) are any two intervals of \(\mathbb{R}\) , each containing more than one point, then \(\mathbf{G}(\mathbf{I})\) and \(\mathbf{G}(\mathbf{J})\) are isomorphic. In \(\mathbf{G}(\mathbf{I})\) , the set \(\mathbf{F}(\mathbf{I})\) of increasing homeomorphisms is a normal subgroup of index 2.
  \end{enumerate}
\end{enumerate}

\begin{enumerate}
\def\labelenumi{\alph{enumi})}
\setcounter{enumi}{1}
\item
  Let \(\mathbf{G} = \mathbf{G}(\mathbf{R})\) and \(\mathbf{F} = \mathbf{F}(\mathbf{R})\) . For each \(f \in \mathbf{G}\) and each \(x \in \mathbf{R}\) , let \(\sigma(x; f)\) denote \(\operatorname{sgn}(f(x) - x)\) . Let \(\mathbf{T}^{+}\) (resp. \(\mathbf{T}^{-}\) ) be the set of all \(f \in \mathbf{F}\) such that \(\sigma(x; f)\) is constant and equal to \(\mathfrak{I}\) (resp. \(-\mathfrak{I}\) ) on \(\mathbf{R}\) . Show that every element of \(\mathbf{T}^{+}\) (resp. \(\mathbf{T}^{-}\) ) is conjugate in \(\mathbf{F}\) to the translation \(x \to x + \mathfrak{I}\) (resp. \(x \to x - \mathfrak{I}\) ). {[}If \(f \in \mathbf{T}^{+}\) , consider the sequence of points \(f^{n}(0), n \in \mathbf{Z}\) {]}.
\item
  Let \(T = T^{+} \cup T^{-}\) . Show that every \(f \in F\) which does not belong to T is the product in F of two elements of \(T_{*}\) {[}Write \(f = f_{1}f_{2}\) in F, where \(f_{1}(x) = \sup (x, f(x))\) and \(f_{2}(x) = \inf (x, f(x))\) ; if g is the translation \(x \to x + 1\) , then \(f_{1}g\) and \(g^{-1}f_{2}\) belong to T{]}.
\item
  Two elements \(f, g\) are conjugate in \(F\) if and only if there exists \(s \in F\) such that \(\sigma(x; f) = \sigma(s(x); g)\) for all \(x \in \mathbb{R}\) . {[}Consider the components \(I_n\) of the open set of all \(x\) such that \(\sigma(x; f) \neq 0\) ; observe that \(F(I_n)\) and \(F\) are isomorphic (by \(a\) ) and use \(b\) {]}. In particular, if \([a, b]\) is an interval of \(R\) such that \(\sigma(a; f) = \sigma(b; f) = 0\) for some \(f \in F\) , then the element \(f'\) of \(F\) which coincides with \(f\) for \(x \leqslant a\) or \(x \geqslant b\) and is such that \(f'(x) = x + \varepsilon(f(x) - x)\) for \(a < x < b\) , where \(0 < \varepsilon < 1\) , is conjugate to \(f\) in \(F\) .
\item
  Let \(f \in F\) and let \(a, b, c, d\) be four real numbers such that \(a < c < b < d\) and \(f(x) - x = 0\) for \(x = a, x = b, x = c, x = d, \sigma(x; f) = 1\) for \(a < x < c\) and \(b < x < d\) . Let \(a', b'\) be two numbers such that \(a \leqslant a' < c\) and \(b < b' \leqslant d\) ; let \(J = [a, b], J' = [a', b']\) , and let \(f_1\) denote the restriction of \(f\) to \(J\) ; thus \(f_1 \in F(J)\) . Show that there is an increasing homeomorphism \(s\) of \(J\) onto \(J'\) such that the element \(f_2 = sf_1s^{-1}\) of \(F(J')\) is such that \(f_2(x) > f^{-1}(x)\) whenever \(a' < x < b'\) {[}use \(d\) {]}.
\end{enumerate}

\(f)\) Let \(f \in \mathbf{F}\) and let \([a, b]\) be an interval of \(\mathbf{R}\) such that \(f(x) - x = 0\) for \(x = a\) and \(x = b\) , \(\sigma(x; f) = -1\) for \(x < a\) and \(\sigma(x; f) = +1\) for \(x > b\) . Show that there exists \(g \in \mathbf{F}\) , conjugate to \(f\) in \(\mathbf{F}\) , such that \(g(x) > f(x)\) for all \(x \in \mathbf{R}\) (same method).

\(g\) ) Let \(\mathbf{H}^{+}\) (resp. \(\mathbf{H}^{-}\) ) denote the normal subgroup of \(\mathbf{F}\) consisting of all \(f\in \mathbf{F}\) such that \(f(x) = x\) in a neighbourhood of \(+\infty\) (resp. \(-\infty\) ). Show that \(\mathbf{H}^{+},\mathbf{H}^{-}\) and \(\mathbf{H} = \mathbf{H}^{+}\cap \mathbf{H}^{-}\) are the only normal subgroups of \(\mathbf{F}\) other than \(\mathbf{F}\) and \(\{e\}\) . {[}If \(\mathbf{N}\) is a normal subgroup of \(\mathbf{F}\) , other than \(\mathbf{F}\) , and if \(f\in \mathbf{N}\) does not belong to \(\mathbf{H}^{+}\) , show that there is an element \(g\in \mathbf{N}\) such that the set of all \(x\in \mathbf{R}\) for which \(\sigma (x;g) = +1\) is an interval \([a, + \infty ]\) , and that we have \(g(x) = x\) for \(x\leqslant a\) . To do this use the constructions of \(e)\) and \(f)\) , as well as \(b)\) and \(c)\) . To show that \(\mathbf{N}\) contains no subgroup, other than \(\mathbf{H}\) and \(\{e\}\) , which is normal in \(\mathbf{F}\) , consider an element \(f\in \mathbf{H}\) distinct from \(e\) ; let \(a\) and \(b\) be the greatest lower bound and least upper bound (both finite by hypothesis) of the set of \(x\in \mathbf{R}\) such that \(\sigma (x;f)\neq 0\) . Observe then that \(\mathrm{F}([a,b])\) is isomorphic to \(\mathbf{F}\) and use the preceding result{]}.

\begin{enumerate}
\def\labelenumi{\alph{enumi})}
\setcounter{enumi}{7}
\tightlist
\item
  Show that the group H is simple. (Same method as for proving that H contains no non-trivial subgroups normal in F).
\end{enumerate}

\begin{enumerate}
\def\labelenumi{\arabic{enumi})}
\setcounter{enumi}{3}
\item
  Let \(f\) be a lower semi-continuous function on a topological space \(\mathbf{X}\) , and let \(\varphi\) be an increasing lower semi-continuous function on \(f(\mathbf{E}) \subset \overline{\mathbf{R}}\) . Show that \(\varphi \circ f\) is lower semi-continuous on \(\mathbf{X}\) .
\item
  Let \(f\) be a lower semi-continuous function on a topological space \(\mathbf{X}\) . Show that if \(\mathbf{A}\) is any non-empty subset of \(\mathbf{X}\) , then \(\sup f(\overline{\mathbf{A}}) = \sup f(\mathbf{A})\) .
\item
  Let \(f\) be a real-valued function defined on a topological space \(X\) . The number (finite or infinite)
\end{enumerate}

\[
\omega (a; f) = \lim _ {x \to a} \sup f (x) - \lim _ {x \to a} \inf f (x)
\]

is called the oscillation of \(f\) at \(a \in \mathbf{X}\) , whenever the right-hand side is defined.

\begin{enumerate}
\def\labelenumi{\alph{enumi})}
\item
  Show that \(x \to \omega(x; f)\) is upper semi-continuous on the subspace \(A\) of \(X\) where this mapping is defined.
\item
  If \(f\) is finite on \(\mathbf{X}\) , then \(\mathbf{A} = \mathbf{X}\) and for each \(a \in \mathbf{X}\) we have
\end{enumerate}

\[
\omega (a; f) = \lim _ {(x, y) \rightarrow (a, a)} \sup (f (x) - f (y)).
\]

\begin{enumerate}
\def\labelenumi{\alph{enumi})}
\setcounter{enumi}{2}
\tightlist
\item
  Let \(f\) be a lower semi-continuous finite real-valued function on \(\mathbf{X}\) . Show that if, at a point \(a \in \mathbf{X}, \omega(a; f)\) has a finite value, then
\end{enumerate}

\[
\liminf _ {x \to a} \omega (x; f) = 0.
\]

{[}Assuming the result false, show that there would exist points x arbitrarily near to a such that \(f(x)\) is as large as we please{]}.

\begin{enumerate}
\def\labelenumi{\alph{enumi})}
\setcounter{enumi}{3}
\tightlist
\item
  For each rational number \(r = p / q\) in irreducible form (with \(q > 0\) ), put \(f(r) = q\) . Show that \(f\) is lower semi-continuous on \(\mathbf{Q}\) , but that at each point \(r \in \mathbf{Q}\) we have \(\omega(r; f) = +\infty\) .
\end{enumerate}

¶ 7) Let X be a locally compact space and f a lower semi-continuous function on X.

\begin{enumerate}
\def\labelenumi{\alph{enumi})}
\tightlist
\item
  If \(\limsup_{x\to a}f(x)=+\infty\) for every \(a\in\mathbf{X}\) , show that the set \(\overline{f}(+\infty)\) is dense in \(\mathbf{X}\) {[}for each \(a\in\mathbf{X}\) and each neighbourhood \(\mathbf{V}\) of \(a\) , define a sequence \((\mathbf{U}_{n})\) of relatively compact open sets such that \(\overline{\mathbf{U}}_{n+1}\subset\mathbf{U}_{n}\) and \(f(x)>n\) for all \(x\in\mathbf{U}_{n}\) {]}.
\end{enumerate}

* b) Let \(n \to r_n\) be a bijection of \(\mathbf{N}\) onto the set of rational numbers belonging to {[}0, 1{]}, and let \(\varphi(x) = \mathrm{i}/\sqrt{|x|}\) {[}with \(\varphi(0) = +\infty\) {]}; then the function \(f(x) = \sum_{n=0}^{\infty} 2^{-n} \varphi(x - r_n)\) is lower semi-continuous on {[}0, 1{]}, \(\overline{f}(+\infty)\) is dense in this interval, and so is its complement {[}to prove the last point, observe that the series is convergent whose general term is

\[
2 ^ {- n} \int_ {0} ^ {1} \varphi (x - r _ {n}) d x. *
\]

(Integration, Chapter IV, § 3, no. 6, Theorem 5).{]} *

\begin{enumerate}
\def\labelenumi{\alph{enumi})}
\setcounter{enumi}{2}
\item
  If \(f\) is finite, show that the set of points \(x\) such that \(\omega(x; f)\) is finite is dense {[}use a){]}.
\item
  Show that the set of points of X at which f is continuous (f being finite or not) is dense. {[}Reduce to the case in which f is bounded, by replacing f by ( f/(1 + \textbar f\textbar) ) (Exercise 4); observe that \(1/\omega(x; f)\) is lower semi-continuous and use a) and Exercise 6c){]}.
\end{enumerate}

\begin{enumerate}
\def\labelenumi{\arabic{enumi})}
\setcounter{enumi}{7}
\item
  Let X be a topological space, and let A be a closed subset of \(X \times \overline{R}\) . Show that the mapping \(x \to \inf(A(x))\) of \(pr_{1}A\) into \(\overline{R}\) is lower semi-continuous. Conversely, if \(f: X \to \overline{R}\) is a lower semi-continuous function, then the subset B of \(X \times \overline{R}\) consisting of pairs \((x, y)\) such that \(y \geqslant f(x)\) is closed in \(X \times \overline{R}\) .
\item
  \begin{enumerate}
  \def\labelenumii{\alph{enumii})}
  \tightlist
  \item
    Let X and Y be two Hausdorff spaces and let \(\pi: X \to Y\) be a proper mapping (Chapter I, § 10, no. 1). Let \(g\) be a lower semi-continuous real-valued function on X. For each \(y \in Y\) , let \(f(y)\) be the greatest lower bound of \(g\) in the set \(\overline{\pi}^{1}(y)\) {[}thus \(f(y) = +\infty\) if \(\overline{\pi}^{1}(y) = \emptyset\) {]}. Show that \(f\) is lower semi-continuous on Y. {[}Use the fact that \(\pi(X)\) is closed, that \(\overline{\pi}^{1}(y)\) is compact for all \(y \in \pi(X)\) , and that every neighbourhood of \(\overline{\pi}^{1}(y)\) contains a neighbourhood which is saturated with respect to the equivalence relation \(\pi(x) = \pi(x')\) .{]}
  \end{enumerate}
\end{enumerate}

\begin{enumerate}
\def\labelenumi{\alph{enumi})}
\setcounter{enumi}{1}
\tightlist
\item
  Let \(g\) be the continuous function \(|x_1 x_2 - 1|\) defined on \(\mathbf{R} \times ]0, +\infty[\) , and for each \(x_1 \in \mathbb{R}\) let \(f(x_1) = \inf_{x_1 > 0} g(x_1, x_2)\) . Show that \(f\) is not lower semi-continuous.
\end{enumerate}

\begin{enumerate}
\def\labelenumi{\arabic{enumi})}
\setcounter{enumi}{9}
\tightlist
\item
  Extend Definition 1, Proposition 1 and Theorem 3 to functions defined on a topological space X which take their values in an arbitrary linearly ordered set Y. If f and g are two mappings of X into Y, both lower semi-continuous at a point a, then inf (f, g) and sup (f, g) are lower semi-continuous at a. So is \(f + g\) if Y is a linearly ordered commutative group.
\end{enumerate}

If \(\mathbf{Y}\) is compact in the topology \(\mathfrak{G}_{0}(\mathbf{Y})\) ( \(\S 2\) , Exercise 6), extend Theorem 4 and Propositions 3 and 4 to functions with values in \(\mathbf{Y}\) .

¶ 11) Let \(\overline{\mathbf{R}}\) carry the right topology (Chapter I, § 1, Exercise 2). A mapping \(f\) of a topological space \(\mathbf{X}\) into \(\overline{\mathbf{R}}\) is continuous in this topology if and only if \(f\) has a relative minimum at every point of \(\mathbf{X}\) . If \(\mathbf{X} = \mathbf{R}\) (with the usual topology), show that the set \(f(\mathbf{X})\) is then countable. {[}Suppose that the result is false, and that the interval \([\alpha, \beta] \cap f(\mathbf{X})\) is uncountable; for each \(y \in [\alpha, \beta]\) , let \(\mathbf{U}(y) = \overline{f}([y, +\infty])\) ; \(\mathbf{U}(y)\) is an open subset in \(\mathbf{R}\) . Let \((\mathbf{I}_n)\) be the sequence (finite or infinite) of components of \(\mathbf{U}(\beta)\) , and for each \(n\) let \(x_n\) be a point of \(\mathbf{I}_n\) ; for each \(y \in [\alpha, \beta]\) , let \(g_n(y)\) (resp. \(h_n(y) \) ) be the left-hand (resp. right-hand) end-point of the component of \(\mathbf{U}(y)\) which contains \(x_n\) . Show that for at least one value of \(n\) one of the functions \(g_n, h_n\) must take an uncountable infinity of distinct values in \([\alpha, \beta]\) ; now use Exercise 8 of § 5{]}.

\begin{enumerate}
\def\labelenumi{\arabic{enumi})}
\setcounter{enumi}{11}
\item
  Let \(f\) be a continuous non-constant finite real-valued function on a compact interval \([a, b]\) of \(\mathbf{R}\) , such that \(f(a) = f(b) = 0\) ; let \(Z\) be the closed subset of \([a, b]\) consisting of points \(x\) for which \(f(x) = 0\) , and let \(c\) be the largest of the lengths of the intervals contiguous to \(Z\) in \([a, b]\) . Show that, for each \(t\) such that \(0 < t < c\) , there is a point \(x \in [a, b]\) such that \(x + t \in [a, b]\) and \(f(x + t) = f(x)\) .
\item
  Let \(f\) be a continuous finite real-valued function on a compact interval \([a, b]\) of \(\mathbf{R}\) , and suppose that \(f\) is not strictly monotone. Show that there exists a point \(x_0 \in ]a, b[\) such that, for each \(\varepsilon > 0\) , there are two points \(y, z\) in \([a, b]\) such that \(x_0 - \varepsilon < y < x_0 < z < x_0 + \varepsilon\) and \(f(y) = f(z)\) .
\item
  Let \(f\) be a continuous mapping of \(\mathbf{R}\) into itself.
\end{enumerate}

\begin{enumerate}
\def\labelenumi{\alph{enumi})}
\tightlist
\item
  Show that if \(f\) is uniformly continuous on \(\mathbf{R}\) , then there exist two real numbers \(\alpha \geqslant 0\) and \(\beta \geqslant 0\) such that \(|f(x)| \leqslant \alpha |x| + \beta\) for all \(x \in \mathbf{R}\) . b) Show that if \(f\) is monotone and bounded in \(\mathbf{R}\) , then \(f\) is uniformly continuous on \(\mathbf{R}\) .
\end{enumerate}

¶ 15) Let X be the Alexandroff half-line {[}§ 2, Exercise 12 b){]}, and let X' be its compactification by adjunction of a point at infinity ω. Show that, for every continuous mapping f: X → X, either we have

\(\lim_{x \to \omega, x \in \mathbf{X}} f(x) = \omega,\) or else there exists \(z \in \mathbf{X}\) such that \(f\) is constant

in the interval \([z, \rightarrow[\) . {[}Show that if, as \(x\) tends to \(\omega, f\) has a cluster point \(c \in \mathbf{X}\) , then we must have \(\lim_{x \to \omega, x \in \mathbf{X}} = c\) ; prove that \(f\) can have no other cluster point in X', by using the fact that every increasing sequence converges in X. Then use Exercise 12 b) of § 2 and the fact that in X' every countable intersection of neighbourhoods of ω is a neighbourhood of ω{]}.

